\section{Establishing the Sim2Real Gap}
\label{sec:gap}

This section addresses \textbf{RQ1}. It quantifies the Sim2Real gap for runway
keypoint detection, identifies which synthetic source generalizes best to real data
overall and per weather scenario, and --- as a central methodological finding ---
shows that a fair measurement of the gap requires controlling for an
\emph{orientation} bias between the synthetic and real test sets, not only for runway
size. The analysis proceeds from a first, counter-intuitive observation (the real set
appears \emph{easier} than the synthetic sets) to its resolution and to a
weather-stratified view that motivates the remainder of the thesis.

\subsection{Experimental Design}
\label{subsec:gap-design}
Seven models are fine-tuned under the full budget (200 epochs, patience 25;
\Cref{subsec:setup-protocol}): one per individual synthetic source (\texttt{flsim},
\texttt{arcgis}, \texttt{bingmaps}, \texttt{ges}, \texttt{xplane}), the
\texttt{combined}/\texttt{all} source, and the extended \texttt{all\_extended}
variant. Each model is evaluated on the test set of every synthetic source and on the
real LARD-V1 set (as a whole and per weather scenario). Since the sources share
\texttt{location\_id}s (\Cref{subsec:setup-data}), the cross-source matrix isolates
the rendering pipeline from the runway viewpoint.

\subsection{Cross-Source Evaluation: A First, Counter-Intuitive Result}
\label{subsec:gap-size}
\Cref{tab:gap-size} reports the cross-source matrix when the test sets are matched on
runway size only. Two patterns emerge. First, among the \emph{synthetic} test sets,
each model performs best on its own source (the diagonal dominates), the expected
signature of a domain effect. Second, and counter-intuitively, the \emph{real} test
set yields the \emph{best} score for almost every training source --- e.g.\
\texttt{ges} reaches $0.653$ on real versus $0.579$ on its own test set. Taken at face
value this \emph{contradicts} the Sim2Real hypothesis, since the models generalize to
real data better than to held-out synthetic data.

\begin{table}[th]
    \caption{Cross-source evaluation, mAP@[.50:.95], test sets matched on runway
    \emph{size} only. Rows: training source; columns: evaluation test set. The real
    column is the best or near-best entry in every row.}
    \label{tab:gap-size}
    \centering
    \begin{NiceTabular}{lrrrrrr}
        \CodeBefore
        \rowcolors{2}{gray!10}{white}
        \Body
        \toprule
        \textbf{Train \textbackslash{} Test} & \textbf{flsim} & \textbf{arcgis} & \textbf{bingmaps} & \textbf{ges} & \textbf{xplane} & \textbf{real} \\
        \midrule
        flsim    & 0.634 & 0.387 & 0.393 & 0.424 & 0.434 & 0.580 \\
        arcgis   & 0.438 & 0.558 & 0.533 & 0.454 & 0.436 & 0.610 \\
        bingmaps & 0.434 & 0.564 & 0.541 & 0.484 & 0.409 & 0.617 \\
        ges      & 0.498 & 0.450 & 0.435 & 0.579 & 0.468 & 0.653 \\
        xplane   & 0.400 & 0.389 & 0.402 & 0.417 & 0.520 & 0.555 \\
        \bottomrule
    \end{NiceTabular}
\end{table}

\subsection{Resolving the Contradiction: The Vertical-Angle Bias}
\label{subsec:gap-orientation}
Investigating the cause revealed a distribution mismatch that the size-only filtering
does not remove: runways in the synthetic test sets frequently appear at oblique
orientations, whereas the real approaches are mostly aligned vertically. Because
off-axis runways are harder, the synthetic sets were effectively \emph{harder} than
the real set, inflating the apparent real performance. Following the definition in
\Cref{subsec:setup-filtering}, the test sets are re-filtered to match \emph{both} the
runway-size and the vertical-angle distributions.

\Cref{tab:gap-angle} reports the resulting matrix. After matching orientation, the
real set is \emph{no longer} the best column: each model now generalizes best to its
own synthetic source (e.g.\ \texttt{flsim} $0.673$ on itself versus $0.557$ on real;
\texttt{ges} $0.644$ versus $0.638$ on real). This \emph{enables} the Sim2Real
hypothesis --- the models do generalize worse to real data than to their training
source. However, the differences are small, so this is best described as a
\emph{modest} rather than a dramatic gap.

\begin{table}[th]
    \caption{Cross-source evaluation, mAP@[.50:.95], test sets matched on runway
    \emph{size and vertical angle}. After controlling for orientation, the diagonal
    (in-source) entry is the best in every row and the real column is no longer the
    maximum.}
    \label{tab:gap-angle}
    \centering
    \begin{NiceTabular}{lrrrrrr}
        \CodeBefore
        \rowcolors{2}{gray!10}{white}
        \Body
        \toprule
        \textbf{Train \textbackslash{} Test} & \textbf{flsim} & \textbf{arcgis} & \textbf{bingmaps} & \textbf{ges} & \textbf{xplane} & \textbf{real} \\
        \midrule
        flsim    & 0.673 & 0.416 & 0.441 & 0.491 & 0.480 & 0.557 \\
        arcgis   & 0.497 & 0.626 & 0.601 & 0.517 & 0.478 & 0.593 \\
        bingmaps & 0.483 & 0.622 & 0.612 & 0.537 & 0.449 & 0.599 \\
        ges      & 0.565 & 0.496 & 0.480 & 0.644 & 0.504 & 0.638 \\
        xplane   & 0.432 & 0.450 & 0.467 & 0.471 & 0.557 & 0.543 \\
        \bottomrule
    \end{NiceTabular}
\end{table}

\subsection{Runway-Size (Distance) Analysis}
\label{subsec:gap-distance}
Breaking performance down by apparent runway size clarifies \emph{where} the gap
lives. For a model trained on \texttt{flsim} and evaluated on its own source versus on
real data, the two curves are nearly equal for small (distant) runways and diverge in
favour of the synthetic source only for larger (nearer) runways. Adding the
vertical-angle criterion shifts the cross-over point from a runway size of $-2.87$ to
$-3.37$ (a difference of $0.5$ on the $\log_{10}$ scale) and shrinks the difference at
larger runways: the gap thus becomes \emph{larger} and is felt at \emph{smaller}
runways than before. For the very smallest runways, the real set still occasionally
outperforms the synthetic source. The same pattern holds when other synthetic sources
replace \texttt{flsim}.

\plotslot{Runway-size vs.\ IoU, model trained on \texttt{flsim}, evaluated on
\texttt{flsim} vs.\ real, before and after adding the vertical-angle filter
(cross-over at $-2.87$ vs.\ $-3.37$).}

\subsection{Why Synthetic Data Is Harder: Image-Quality Causes}
\label{subsec:gap-quality}
That sim-trained models generalize surprisingly well --- and sometimes better --- to
real data than to their own synthetic source, especially relative to \emph{other}
synthetic sources, points to a difference in \emph{image quality} rather than a small
domain gap. The effect is most pronounced for distant runways and makes runway
detection on synthetic data genuinely harder. Comparing synthetic and real images at
matched runway size reveals three recurring causes:
\begin{itemize}
    \item \textbf{Runway/background separability.} On synthetic imagery, distant
    runways are often hard to distinguish from their pixel surroundings, leading to
    missed or poorly localized detections relative to equally distant real runways
    (which are rendered/captured at higher effective resolution).
    \item \textbf{Confusion with other instances.} Runways are frequently confused
    with parallel taxiways/roads or with neighbouring runways. This happens on both
    domains but appears more often on synthetic data.
    \item \textbf{Edge sharpness.} Runway edges tend to be sharper in real images than
    in synthetic ones, which yields higher precision and, almost always, a higher IoU.
\end{itemize}
A further, systematic error is independent of image quality: in \emph{near}
perspectives the model consistently under-covers the runway threshold markings (the
``piano keys''), predicting only the near portion of the quadrilateral. Other
mismatches --- time of day, lighting contrast and runway appearance --- also
contribute to isolated failures.

\plotslot{Qualitative examples at matched runway size: synthetic vs.\ real, ground
truth (green) and YOLOv8-pose prediction (blue), with per-image Prec/Rec/IoU
(separability, instance confusion, edge sharpness).}

\subsection{Weather-Stratified Analysis}
\label{subsec:gap-weather}
Splitting the (size- and angle-filtered) real set into weather scenarios shows that
the overall real performance is carried by the dominant nominal split and hides a
large spread across conditions. \Cref{tab:gap-weather} reports mAP@[.50:.95] per
scenario. Three observations stand out:
\begin{itemize}
    \item \textbf{Performance depends strongly on the scenario}, and \emph{fog} is by
    far the hardest ($\approx 0.21$--$0.32$ across sources).
    \item \textbf{The best source is scenario-dependent.} The satellite-imagery
    sources (\texttt{ges}, \texttt{arcgis}, \texttt{bingmaps}) are strongest under
    \emph{glare} and \emph{nominal}, whereas the full 3D simulator \texttt{flsim} is
    strongest under \emph{rain}, \emph{fog} and \emph{night}.
    \item \textbf{The combined source generalizes best overall.} \texttt{all}/\texttt{all\_extended}
    is best or near-best in almost every scenario (e.g.\ nominal $0.636$, night
    $0.524$), confirming that \emph{rendering-source diversity outweighs
    runway-viewpoint coverage} in this setup, even though \texttt{combined} covers
    only about one fifth as many locations.
\end{itemize}
The good aggregate real performance is therefore largely attributable to the nominal
split's dominance; the adverse splits reveal the domain shift that the remaining
techniques target.

\begin{table}[th]
    \caption{Real test set per weather scenario, mAP@[.50:.95] (size- and
    angle-filtered), by training source. \texttt{combined} = \texttt{all}. Fog is the
    hardest scenario; the best source varies by scenario.}
    \label{tab:gap-weather}
    \centering
    \begin{NiceTabular}{lrrrrr}
        \CodeBefore
        \rowcolors{2}{gray!10}{white}
        \Body
        \toprule
        \textbf{Train source} & \textbf{Nominal} & \textbf{Glare} & \textbf{Rain} & \textbf{Fog} & \textbf{Night} \\
        \midrule
        flsim         & 0.533 & 0.443 & 0.596 & 0.315 & 0.376 \\
        arcgis        & 0.586 & 0.507 & 0.400 & 0.224 & 0.268 \\
        bingmaps      & 0.589 & 0.539 & 0.425 & 0.209 & 0.278 \\
        ges           & 0.628 & 0.538 & 0.460 & 0.210 & 0.262 \\
        xplane        & 0.533 & 0.450 & 0.286 & 0.221 & 0.282 \\
        all           & 0.636 & 0.586 & 0.527 & 0.280 & 0.406 \\
        all\_extended & 0.623 & 0.570 & 0.603 & 0.315 & 0.524 \\
        \bottomrule
    \end{NiceTabular}
\end{table}

To understand \emph{why} different sources suit different scenarios, the datasets are
characterized in a style-oriented feature space, using (i) a small interpretable
descriptor (brightness, contrast, saturation, sharpness as the variance of the
Laplacian, edge density, colorfulness) and (ii) a texture descriptor given by the
averaged Gram matrices of a VGG16 network pretrained on ImageNet-1k. A UMAP projection
of these features separates the sources into clusters and places the real domain apart
from all synthetic sources, confirming a genuine style gap and motivating the
appearance-level interventions of \Cref{sec:weather-aug,sec:style-transfer}.

\plotslot{UMAP of style-oriented image features (interpretable + VGG-Gram) for the
five synthetic sources and the real domain, showing a distinct ``real'' cluster.}

\subsection{Discussion}
\label{subsec:gap-discussion}
The gap is real but modest once orientation is controlled for, and the surprisingly
strong real-domain performance is largely explained by higher real image quality
rather than by the absence of a domain shift. Two takeaways are worth carrying
forward. First, cross-domain comparisons for runway detection \emph{must} control for runway
size \emph{and} orientation, or they will systematically flatter the real set.
Second, the gap concentrates in adverse weather --- especially fog, rain and night ---
and no single synthetic source dominates all scenarios, although the combined source
comes closest. From here on, the techniques target \textbf{generalization to real
adverse-weather conditions} when training on nominal-condition synthetic data.
